\documentclass[11pt]{scrartcl}
\usepackage{geometry}
\usepackage{hyperref}
\usepackage{xcolor}
\usepackage{amsfonts}
\usepackage{pifont}
\usepackage{enumitem}
\usepackage{minted}
\usepackage[sfdefault]{carlito}
\renewcommand{\familydefault}{\sfdefault}
\geometry{margin=1in}

\definecolor{xhbg}{HTML}{FAFAFA}
\setminted{
  fontsize=\small,
  breaklines=true,
  autogobble=true,
  frame=single,
  framesep=2mm,
  linenos=true,
  tabsize=2,
  bgcolor=xhbg
}
\usemintedstyle{friendly}

\newminted[code]{haskell}{}
\newcommand{\tododone}{\textcolor{green!45!black}{\ding{51}}}
\newcommand{\todopending}{\textcolor{red!65!black}{\ding{55}}}
\newcommand{\todoplanned}{\textcolor{blue!70!black}{$\mathbb{P}$}}

\title{Xeus-Haskell: Docs and Literate Sources}
\author{Masaya Taniguchi and xeus-haskell contributors}
\date{}

\begin{document}
\maketitle
\tableofcontents
\newpage

\section{Welcome to Xeus-Haskell}


Xeus-Haskell is a Jupyter kernel for Haskell based on the native implementation of the Jupyter protocol, \href{https://github.com/jupyter-xeus/xeus}{xeus}.
\subsection{Preliminaries}


To start using or developing Xeus-Haskell, please note the following:
\begin{itemize}
\item \textbf{MicroHs Backend:} Xeus-Haskell uses \href{https://github.com/augustss/MicroHs}{MicroHs}, a minimal implementation of Haskell. Therefore, supported libraries and extensions are limited.
\item \textbf{Native Protocol:} Based on \href{https://github.com/jupyter-xeus/xeus}{xeus} (C++), so the kernel must be compiled for the specific platform you are using.
\item \textbf{WebAssembly Support:} You can compile this kernel to WebAssembly to run it directly in browsers!
\end{itemize}


\subsection{Quickstart}


Since this project is not yet available on Conda-forge or Emscripten-forge, you need to build it from source. We have organized the workflows using \href{https://pixi.sh}{pixi}.
\subsubsection{For Native JupyterLab}


\begin{minted}{text}
git clone https://github.com/jupyter-xeus/xeus-haskell
cd xeus-haskell
pixi run -e default prebuild
pixi run -e default build
pixi run -e default install
pixi run -e default serve # JupyterLab is ready!
\end{minted}


\subsubsection{For WebAssembly JupyterLite}


\begin{minted}{text}
git clone https://github.com/jupyter-xeus/xeus-haskell
cd xeus-haskell
pixi install -e prebuild
pixi run -e wasm-build prebuild
pixi run -e wasm-build build
pixi run -e wasm-build install
# pixi run -e wasm-build fix-emscripten-links # Run this if you encounter linking errors
pixi run -e wasm-build serve # JupyterLite is ready!
\end{minted}


\subsection{A Deep Dive into Xeus-Haskell}


\begin{itemize}
\item \hyperref[sec:contributing]{Contributing Guide}
\item \hyperref[sec:communication]{Communication Protocol}
\item \hyperref[sec:repl]{REPL Mechanism}
\item \hyperref[sec:todo]{Roadmap \& TODO}
\item \href{https://github.com/jupyter-xeus/xeus-haskell/blob/main/AGENTS.md}{AI Agent Guide}
\end{itemize}

\section{Communication and System Architecture}
\label{sec:communication}


This document describes the internal communication flow of \texttt{xeus-haskell}, from the Jupyter frontend down to the MicroHs runtime.
\subsection{System Stack}


The communication stack is organized as follows:
Jupyter Frontend \(\leftrightarrow\) Xeus C++ Library \(\leftrightarrow\) Xeus-Haskell Kernel \(\leftrightarrow\) \texttt{mhs\_repl.cpp} Bridge \(\leftrightarrow\) Repl.lhs Compiler API \(\leftrightarrow\) MicroHs Runtime.


\subsection{Layer Details}


\subsubsection{Jupyter and Xeus Protocol Layer}


This layer handles the ZeroMQ-based Jupyter protocol. It manages the heartbeat, shell, and control sockets.
\subsubsection{Xeus to Kernel Interpreter Layer}


The \texttt{xinterpreter} implementation hooks into the Xeus kernel lifecycle. It receives execution requests and delegates them to the Haskell backend.
\subsubsection{C++ Bridge to Haskell Runtime Layer}


This is the bridge layer between C++ and the Haskell-compiled Compiler API wrapper (\texttt{Repl.lhs}).
The \texttt{mhs\_repl.cpp} bridge performs three primary tasks:
\begin{enumerate}
\item \textbf{Initialize MicroHs Runtime}:
   The \texttt{MicroHsRepl} constructor locates runtime files (standard libraries) using \texttt{microhs\_runtime\_dir}. After initialization, it evaluates a trivial expression (\texttt{"0"}) to warm up the compiler and cache library compilation.
\item \textbf{Capture Standard Output}:
   To receive output from the Haskell world, we use POSIX pipes to redirect \texttt{stdout}.
   \begin{itemize}
   \item C++ sets up a redirection from \texttt{stdout} to a temporary buffer.
   \item After Haskell evaluation completes, the buffer is read.
   \item For implementation details, see \texttt{capture\_stdout} in the source code.
   \end{itemize}
\item \textbf{Parse Captured Content}:
   The output is parsed to distinguish plain text from rich content (HTML, LaTeX, etc.) using the internal protocol described below.
\end{enumerate}


\subsubsection{4. REPL Engine and MicroHs Compiler Layer}


This layer uses the MicroHs Compiler API to evaluate Haskell code strings in the runtime environment.
\subsection{Rich Content Protocol}


\texttt{xeus-haskell} supports rendering rich media (HTML, images, LaTeX) by embedding control sequences in the standard output stream.
\subsubsection{Protocol Structure}


Rich content is strictly delimited by \textbf{ASCII control characters}:
\begin{minted}{text}
<STX><MIME_TYPE><US><CONTENT><ETX>
\end{minted}


\begin{center}
\begin{tabular}{|l|c|p{8cm}|}
\hline
\textbf{Component} & \textbf{ASCII (Hex)} & \textbf{Description} \\
\hline
\texttt{STX} & \texttt{02} & \textbf{S}tart of \textbf{T}e\textbf{x}t \\
\hline
\texttt{MIME\_TYPE} & - & The target MIME type (e.g., \texttt{text/html}) \\
\hline
\texttt{US} & \texttt{1F} & \textbf{U}nit \textbf{S}eparator \\
\hline
\texttt{CONTENT} & - & The raw data/string to be rendered \\
\hline
\texttt{ETX} & \texttt{03} & \textbf{E}nd of \textbf{T}e\textbf{x}t \\
\hline
\end{tabular}
\end{center}
\subsubsection{Supported Content Types}


\begin{itemize}
\item \texttt{text/plain}: Standard console output.
\item \texttt{text/html}: Rendered as HTML in Jupyter.
\item \texttt{text/latex}: Mathematical equations rendered via MathJax.
\item \texttt{text/markdown}: Formatted text and tables.
\end{itemize}


All content is expected to be \textbf{UTF-8} encoded.
\subsection{Haskell Display API}


The protocol is abstracted in Haskell through the \texttt{Display} typeclass.
\subsubsection{Display Typeclass Interface}


To enable rich output for a custom type, implement the \texttt{Display} typeclass:
\begin{minted}{haskell}
import XHaskell.Display

-- | Data structure for the protocol
data DisplayData = DisplayData {
    mimeType :: String,
    content  :: String
}
\end{minted}


\subsubsection{Implementation Examples}


\begin{minted}{haskell}
newtype HTMLString = HTMLString String
instance Display HTMLString where
    display (HTMLString s) = DisplayData "text/html" s

newtype LaTeXString = LaTeXString String
instance Display LaTeXString where
    display (LaTeXString s) = DisplayData "text/latex" s

-- Default instance for plain strings
instance Display String where
    display s = DisplayData "text/plain" s
\end{minted}


\subsubsection{Protocol Generation}


The \texttt{DisplayData} record implements a \texttt{Show} instance that automatically formats the string into the \texttt{<STX>...<ETX>} protocol.
\begin{minted}{haskell}
putStr $ show (display (HTMLString "<p style='color: red'>Hello Haskell!</p>"))
\end{minted}


\texttt{02} \texttt{text/html} \texttt{1F} \texttt{<p style='color: red'>Hello Haskell!</p>} \texttt{03}

\section{Contributing to xeus-haskell}
\label{sec:contributing}


Thank you for your interest in contributing to \texttt{xeus-haskell}!
Xeus and xeus-haskell are subprojects of Project Jupyter and are subject to the \href{https://github.com/jupyter/governance}{Jupyter governance} and \href{https://github.com/jupyter/governance/blob/master/conduct/code\_of\_conduct.md}{Code of conduct}.
\subsection{General Guidelines}


For general documentation about contributing to Jupyter projects, see the \href{https://jupyter.readthedocs.io/en/latest/contributor/content-contributor.html}{Project Jupyter Contributor Documentation}.
\subsection{Community}


The Xeus team organizes public video meetings. The schedule for future meetings and minutes of past meetings can be found on our \href{https://jupyter-xeus.github.io/}{Team Compass}.
\subsection{Getting Started}


\subsubsection{Prerequisites}


The project uses \href{https://pixi.sh/}{Pixi} for dependency management and workflow automation. If you haven't installed it yet:
\begin{minted}{text}
curl -fsSL https://pixi.sh/install.sh | sh
\end{minted}


\subsubsection{Setting Up Your Environment}


\begin{enumerate}
\item Fork the repository on GitHub.
\item Clone your fork locally:
\end{enumerate}


\begin{minted}{text}
git clone https://github.com/<your-username>/xeus-haskell.git
cd xeus-haskell
\end{minted}


\begin{enumerate}
\item Initialize the development environment:
\end{enumerate}


\begin{minted}{text}
pixi run -e default prebuild
\end{minted}


The \texttt{default} environment includes all necessary tools (CMake, C++ compilers, Python for testing, etc.).
\subsection{Development Workflow}


\subsubsection{Building the Kernel}


To build and install the kernel in your local environment:
\begin{minted}{text}
pixi run -e default build
pixi run -e default install
\end{minted}


\subsubsection{Running Tests}


We use \texttt{pytest} and \texttt{ctest} for verifying the kernel:
\begin{minted}{text}
# Run C++ tests
pixi run -e default ctest

# Run Python-based Jupyter kernel tests
pixi run -e default pytest
\end{minted}


\subsubsection{Working with WebAssembly (JupyterLite)}


If you are contributing to the WebAssembly build, you need the \texttt{wasm-build} and \texttt{wasm-host} environments:
\begin{minted}{text}
# Prepare the WASM host environment
pixi install -e wasm-build
pixi install -e wasm-host

# Build for WebAssembly
pixi run -e wasm-build prebuild
pixi run -e wasm-build build
pixi run -e wasm-build install

# Run the JupyterLite server locally
pixi run -e wasm-build serve
\end{minted}


\subsection{Project Structure}


\begin{itemize}
\item \texttt{src/}: C++ source code for the kernel and REPL bridge.
\item \texttt{include/}: C++ header files.
\item \texttt{share/}: Haskell side of the implementation, including \texttt{Repl.lhs} and MicroHs libraries.
\item \texttt{test/}: Python tests for kernel functionality.
\item \texttt{xhaskell.tex}: Integrated LaTeX manual source that includes all architecture chapters and literate Haskell modules.
\end{itemize}


\subsection{Style Guidelines}


\begin{itemize}
\item \textbf{C++}: Follow modern C++ practices. We use Xeus as the base library.
\item \textbf{Haskell}: Since we use MicroHs, avoid dependencies or language features not supported by the MicroHs compiler.
\item \textbf{Documentation}: Write documentation in LaTeX (\texttt{.tex}) and keep examples runnable with the current Pixi workflows.
\end{itemize}

\subsubsection{Building Documentation}

Generate the project PDF documentation with:
\begin{minted}{text}
pixi run -e docs generate
\end{minted}


\subsection{Submitting Changes}


\begin{enumerate}
\item Create a new branch for your feature or bugfix.
\item Ensure all tests pass.
\item Submit a Pull Request with a clear description of the changes.
\end{enumerate}


\subsection{Community}


Join our public video meetings! The schedule and minutes are available on our \href{https://jupyter-xeus.github.io/}{Team Compass}.

\section{TODO: Unimplemented Features \& Roadmap}
\label{sec:todo}


This document tracks planned improvements and currently missing features for \texttt{xeus-haskell}.

Status symbols: \tododone\ completed, \todopending\ pending, \todoplanned\ planned.
\subsection{Core Kernel Features}


\begin{itemize}[leftmargin=*,itemsep=0.45em]
\item[\tododone] \textbf{\texttt{is\_complete} implementation}: Map parsability of mixed cells to Jupyter's completeness status.
\item[\todopending] \textbf{Incomplete Status}: Refine \texttt{is\_complete} to distinguish between \texttt{invalid} and \texttt{incomplete} code (e.g., unclosed delimiters).
\item[\tododone] \textbf{\texttt{inspect\_request}}: Support for "Introspection" (Shift+Tab in Jupyter) to show documentation or type signatures for identifiers.
\item[\todopending] \textbf{\texttt{history\_request}}: Implementation of the Jupyter history protocol to allow searching and retrieving previous cell inputs.
\item[\todopending] \textbf{Advanced Completion}: Improve \texttt{completion\_request} to support qualified names (e.g., \texttt{Prelude.putStrLn}) and type-aware suggestions.
\item[\todopending] \textbf{Kernel Interrupt}: Add support for interrupting long-running Haskell executions (SIGINT handling).
\end{itemize}


\subsection{REPL Improvements}


\begin{itemize}[leftmargin=*,itemsep=0.45em]
\item[\todopending] \textbf{Selective Re-compilation}: Instead of concatenating all previous definitions for every new definition, implement a dependency-tracking system to only re-compile what is necessary.
\item[\todopending] \textbf{Better Error Reporting}: Map MicroHs compile errors back to the specific line numbers in the Jupyter cell.
\item[\todopending] \textbf{WASM Optimization}: Reduce the size of the WebAssembly bundle and optimize the initial "warmup" time in the browser.
\item[\tododone] \textbf{Language Extensions}: Work with the MicroHs upstream to support more modern Haskell extensions (e.g., \texttt{GADTs}, \texttt{TypeFamilies}).
\end{itemize}


\subsection{Rich Output Display System}


\begin{itemize}[leftmargin=*,itemsep=0.45em]
\item[\todoplanned] \textbf{Standard Library Instances}: Add \texttt{Display} instances for more types in the MicroHs standard library.
\item[\todoplanned] \textbf{Data Visualization}: Create bridges for common Haskell charting libraries to output Jupyter-compatible JSON/Vega-Lite data.
\item[\todoplanned] \textbf{Widgets}: Initial investigation into supporting Jupyter Widgets (ipywidgets protocol).
\end{itemize}


\subsection{Infrastructure \& Testing}


\begin{itemize}[leftmargin=*,itemsep=0.45em]
\item[\todoplanned] \textbf{Enhanced CI}: Add automated UI tests using \texttt{playwright} or \texttt{selenium} to verify the kernel works in a real JupyterLab/JupyterLite instance.
\item[\todoplanned] \textbf{Benchmarking}: Create a suite of performance benchmarks to track compilation and execution speed regressions.
\item[\todoplanned] \textbf{Conda-forge/Emscripten-forge}: Package the kernel for easier installation without needing to build from source.
\end{itemize}

\clearpage
\section{Literate Haskell Source}
\input{src/Repl.lhs}
\subsection{Snippet Analysis and Parsing Module}
\input{src/Repl/Analysis.lhs}
\subsection{Compilation and Execution Translation Module}
\input{src/Repl/Compiler.lhs}
\subsection{Persistent Context and Definition Storage Module}
\input{src/Repl/Context.lhs}
\subsection{Error Classification and Rendering Module}
\input{src/Repl/Error.lhs}
\subsection{Runtime Command Execution Module}
\input{src/Repl/Executor.lhs}
\subsection{Shared Utility Functions Module}
\input{src/Repl/Utils.lhs}

\end{document}
